\section{One ordinary period for the nonprincipal CM ideal bases}
\label{ci-section}

The source permits a nonprincipal elliptic lattice.  At every
prime where its moment theorem is applied it also assumes
$p\nmid N(I)$.  We verify the extension of
Section~\ref{ct-section} under this precise hypothesis, using a
prime-to-$p$ isogeny to keep the same period.  Choosing an unrelated
unit in $\mathbb C_p$ and making a binomial substitution with that
unit would not justify this step.

Keep $k=\mathbb Q(\sqrt{-D})$ and
$\mathcal O_k=\mathbb Z+\mathbb Z\tau$ as in
Section~\ref{tc-section}.  Throughout this section $p\geq5$
is a rational prime split in $k$; this is the scope of the
ordinary theta integrality theorem used here.
Let $I$ be the ideal (TC43),
so $[\mathcal O_k:I]=a$, $a$ is odd, and $p\nmid a$.
The source's curve ideals $J$ satisfy
$\gcd(N(J),2D)=1$, in particular $N(J)$ is odd.
Such an integral ideal $J$ can
first be written $J=gI$ with $g\in\mathbb Z_{>0}$ and
$I$ primitive of this form.  One obtains $g$ by taking the
content of a lattice basis; stability under $\tau$ shows
that the primitive Hermite normal form has second
$\tau$-coefficient one.
Then $N(J)=g^2a$, so $p\nmid N(J)$ implies $p\nmid a$
and $p\nmid g$.  The factor $g$ is included in the
scaling treated at the end of this section.
Put
\[
 H_I(z,w)=\frac{2}{a\sqrt D}z\overline w,\qquad
 E_I=\operatorname{Im}H_I,\qquad
 Q_I(z)=\frac{|z|^2}{a},\qquad h_I=\frac a2.
 \tag{CI1}
\]
Let $\theta_I$ be the reduced odd theta function translated
by $h_I$ and divided by its value at zero.  Lemma~\ref{tc-ideal-character}
gives its multiplier $(-1)^{Q_I(v)}$ for $v\in I$.
It is even and has its zero precisely at $h_I+I$, by the
same zero-count and uniqueness argument as in
Lemma~\ref{ct-even-identification}.  All translations below use
the Hermitian form $H_I$.

\subsection{Prime-to-$p$ isogenies and the common formal parameter}

Let $\mathscr E_0$ be the good ordinary model for
$\mathbb C/\mathcal O_k$ used in Section~\ref{ct-section}.
We choose its CM equation and differential over algebraic
numbers before completing at the selected place, as is
permitted in that classical CM input.  Further finite
extensions and the quotients below preserve this convention.
Write its period lattice relative to a good differential as
$c_0\mathcal O_k$, and normalize its multiplicative formal
parameter there.  If $\Omega_{p,0}$ is the associated unit
period, put
\[
                  \omega=\frac{\Omega_{p,0}}{\iota(c_0)}.
 \tag{CI2}
\]
The embedding $\iota$ is the fixed one from that section.

\begin{proposition}\label{ci-common-period}
There is a good ordinary model $\mathscr E_I$ for
$\mathbb C/I$ whose multiplicative formal parameter, pulled
back through the isogeny
\[
                    \psi:\mathbb C/I\longrightarrow
                         \mathbb C/\mathcal O_k,\qquad z\longmapsto z,
 \tag{CI3}
\]
has the same formal relation $z=\omega\log(1+T)$
in the unscaled complex variable.
If $\delta_I^{12}=g_2(I)^3-27g_3(I)^2$ and $\delta_0$
is the corresponding root for $\mathcal O_k$, then
\[
 \delta_I/\delta_0\in\overline{\mathbb Q}^{\,\times},
       \qquad |\iota(\delta_I/\delta_0)|_p=1.
 \tag{CI4}
\]
\end{proposition}
\begin{proof}
The inclusions $a\mathcal O_k\subset I\subset\mathcal O_k$
give (CI3) and its degree-$a$ dual, which in unscaled complex
coordinates sends $z$ to $az$.
Construct the latter as a quotient of $\mathscr E_0$ by its
subgroup $a^{-1}I/\mathcal O_k$.
After a finite extension this subgroup extends to a finite
\'etale subgroup of $\mathscr E_0[a]$: since $p\nmid a$,
multiplication by $a$ has invertible differential, and its
kernel on the smooth elliptic scheme is finite \'etale.
The translation action is free.  The finite locally free
quotient theorem \cite{stacks-07S7} makes the quotient a scheme
and its quotient map a torsor.  In this case it is finite
\'etale; smoothness and properness descend through that
cover.  Its geometrically connected genus-one fibers and
the image of the zero section make the quotient an elliptic
scheme.  Thus it is the required good model $\mathscr E_I$.
The dual isogeny (CI3) extends as well.  Both maps are
isomorphisms of formal groups: their composite is $[a]$,
whose formal inverse is integral, and each has a unit
linear term.  Ordinary reduction is preserved, for example
because they induce an isomorphism on the $p$-divisible
formal groups.

Write the period lattice of $\mathscr E_I$, relative to a
good differential, as $c_I I$.
The differential multiplier of $\psi$ in these coordinates
is $u=c_0/c_I$.  It is an algebraic number and a $p$-adic
unit, since $\psi$ is an algebraic isogeny between the
chosen algebraic models and its map on formal groups is an
integral isomorphism.
If $\lambda_I,\lambda_0$ are their normalized formal
logarithms, differentiation of the group homomorphism gives
\[
               \lambda_0(\psi(t))=\iota(u)\lambda_I(t).
 \tag{CI5}
\]
There is no extra constant because both sides vanish at zero.
Pull back the chosen multiplicative coordinate of
$\mathscr E_0$ along $\psi$.  Its period on $\mathscr E_I$
is
$\Omega_{p,I}=\Omega_{p,0}/\iota(u)
=\iota(c_I/c_0)\Omega_{p,0}$, again a unit.
In the variable $z$ for the lattice $I$, it therefore gives
\[
             \frac{\Omega_{p,I}}{\iota(c_I)}
                       =\frac{\Omega_{p,0}}{\iota(c_0)}=\omega.
 \tag{CI6}
\]
This proves the asserted equality of periods, rather than
merely equality of their absolute values.

Let $\Delta_0,\Delta_I$ be the discriminants of the two good
models.  The invariant-differential calculation (CT11) gives
$(\delta_I/c_I)^{12}=\Delta_I$ and
$(\delta_0/c_0)^{12}=\Delta_0$, after the fixed algebraic
identifications.  Their ratio shows that
$(\delta_I/\delta_0)^{12}
=(c_I/c_0)^{12}\Delta_I/\Delta_0$ is algebraic.
All its factors have $p$-adic absolute value one.
Taking a twelfth root proves both assertions in (CI4).
\end{proof}

\begin{theorem}\label{ci-integral-translates}
For every $t\in I\otimes\mathbb Q$ of order prime to $p$
modulo $I$,
\[
                     \widehat{U_t\theta_I}
                           \bigl(\omega\log(1+T)\bigr)
                         \in\mathcal O_{\mathbb C_p}[[T]].
 \tag{CI7}
\]
The series is a unit unless $t+h_I\in I$.
All finite-combination norm identities of
Theorem~\ref{ct-combination-norm} hold with $\theta_I$ and
this same $\omega$.
In particular every prime-to-$p$ translate has integral
period-scaled moments in every finite family of directions
from $I$.
\end{theorem}
\begin{proof}
Apply the reduced-theta integrality theorem cited in
Section~\ref{ct-section} to the good ordinary curve
$\mathscr E_I$ and to the nonzero two-torsion point $h_I$.
Its reduced-theta value at zero is a unit.  Translation
by $t$ composes with that translation up to the root of
unity $e^{-\pi iE_I(t,h_I)}$.
The sum $t+h_I$ still has prime-to-$p$ order.
The cited theorem therefore gives integrality, with unit
constant term except at the stated zero.
The inverse of the multiplicative coordinate constructed
in Proposition~\ref{ci-common-period} is integral with
unit linear coefficient, and (CI6) identifies its variable
with the one in (CI7).

The finite-combination assertion follows from the proof of
Theorem~\ref{ct-combination-norm} on this elliptic scheme:
choose an even common denominator $N$ prime to $p$, use the
line bundle $[N]^*\mathcal O([h_I])^{\otimes m}$, and use
the integral trivialization near zero given by its canonical
section.  Formula (CI7) makes the corresponding reduced
expansion an integral unit.  The saturated expansion lemma
identifies its coefficient norm with the integral
global-section norm.  Translation by $t/N$ has a power
equal to a root of unity and is therefore an integral
automorphism by the vertical-divisor proof of
Lemma~\ref{ct-torsion-isometry}.
The exact identity (CT18) uses only the Hermitian Weyl
formula, so it holds for $H_I$ as well.
Finally $T\mapsto(1+T)^N-1$ and its inverse are integral,
and removing that substitution gives the asserted norm
identity for arbitrary $\mathbb C_p$ coefficients.
Every element of $I$ belongs to $\mathcal O_k$, so its
image at this split prime lies in $\mathbb Z_p$.
Lemma~\ref{lm-directions} now gives the assertion for
arbitrary finite direction families.
\end{proof}

\subsection{The selected torsion character and its Weyl phase}

The chosen parameter is compatible with the required
character, not just with the invariant differential.

\begin{proposition}\label{ci-character}
The formal $p^n$-torsion of $\mathbb C/I$ consists of
\[
 y=\frac{am+b'\tau_b}{p^n}\pmod I,\qquad
 \iota(am+b'\overline{\tau_b})\in p^n\mathbb Z_p.
 \tag{CI8}
\]
In the parameter of Proposition~\ref{ci-common-period},
\[
                1+T(y)=\iota(e^{2\pi ib'/p^n})
                         =\iota(e^{-2\pi iE_I(a,y)}).
 \tag{CI9}
\]
If $st\in I$, $p\nmid s$, and
$y=s\nu/p^n$ with $\nu\in I$ and
$\iota(\bar\nu)\in p^{2n}\mathbb Z_p$, then
\[
 \widehat{e^{\pi iQ_I(y)}U_y(U_t\theta_I)}
                       \bigl(\omega\log(1+T)\bigr)
  =S_\zeta\!\left(\widehat{U_t\theta_I}
                       \bigl(\omega\log(1+T)\bigr)\right),
 \quad
 \zeta=\iota(e^{-2\pi iE_I(a,y)}).
 \tag{CI10}
\]
\end{proposition}
\begin{proof}
The isogeny (CI3) identifies the formal groups and their
$p$-power torsion.  Since
$\tau_b=\tau+(b-1)/2$, the image of the numerator in
(CI8), in the basis $1,\tau$, is
$am+b'(b-1)/2+b'\tau$.
Its conjugate-divisibility condition and its coefficient
of $\tau$ are exactly those in
Proposition~\ref{ct-character-normalization}.
This proves (CI8)--(CI9), including the simultaneous
normalization for every $n$.
Alternatively, $a$ is a unit modulo $p^n$, so (CI8)
has exactly one residue of $m$ for every $b'$; this
also proves directly that its character is an isomorphism.

We record the arithmetic changes in the proof of
Theorem~\ref{ct-power-translation}; they are the only
changes needed for (CI10).
Because $N(\nu)/a\in\mathbb Z$ and $p\nmid a$, the strong
conjugate divisibility gives
$Q_I(\nu/p^n)\in\mathbb Z$.
Choose $\epsilon\in2\mathcal O_k$ with
$\epsilon\equiv1\pmod{\mathfrak p^n}$ and
$\epsilon\equiv0\pmod{\bar{\mathfrak p}^{\,n}}$.
Then
\[
  (\epsilon-1)\nu/p^n\in I,\qquad
                  \bar\epsilon\nu/p^n\in2I.
 \tag{CI11}
\]
At the two primes over $p$ this follows from the two
congruences and the strong divisibility, using
$I\otimes\mathbb Z_p=\mathcal O_k\otimes\mathbb Z_p$.
At every other prime the denominator is a unit;
the factor two follows because $\epsilon$ is even and
$p$ is odd.
The rational elliptic function
$\theta_I(\epsilon z)/\theta_I(\bar\epsilon z)$ has no
zero or pole in the formal subgroup: all such points
map to the nonzero two-torsion point $h_I$.
Applying the formal translation identity (CT22) and
cancelling its unit denominator, exactly as in
(CT28)--(CT30), gives
$S_\zeta\widehat\theta_I
=\widehat{U_{\epsilon y}\theta_I}$ for $y=\nu/p^n$.

Write $\epsilon=A+B\tau$ and $l=(\epsilon-1)y$.
The two phase identities are now
\[
 E_I(y,l)=-B Q_I(y),\qquad
 Q_I(l)=(N(\epsilon)-2A-B+1)Q_I(y).
 \tag{CI12}
\]
Consequently
$e^{\pi iE_I(y,l)-\pi iQ_I(l)}
=e^{\pi iQ_I(y)}$, since their exponent difference
is $(2A-2-N(\epsilon))Q_I(y)$, an even integer.
This proves the base-theta assertion.
For $st\in I$, apply (CT22) to the elliptic rational
function
$z\mapsto(U_t\theta_I)(sz)/\theta_I(sz)$ and to
$\nu/p^n$.  Its poles avoid the formal subgroup.
The invertible substitution $[s]$ gives the ratio
identity (CT33), now for $I$.
Multiplying by the base-theta identity proves (CI10).
The operator $S_\zeta$ and its inverse preserve the
coefficient norm by (CT21).
\end{proof}

\subsection{Scaling the lattice and the actual restriction on primes}

For $L=\lambda I$ and its normalized theta function,
uniqueness gives $\theta_L(\lambda z)=\theta_I(z)$, and
the Hermitian forms give
\[
          (U_{\lambda t}^{L}\theta_L)(\lambda z)
                           =(U_t^{I}\theta_I)(z).
 \tag{CI13}
\]
Thus a direction $v_j=\lambda u_j\in L$ satisfies, identically
before transport,
\[
 (U_{\lambda t}^{L}\theta_L)
       \left(\omega\sum_j\lambda u_j\log(1+T_j)\right)
 =
 (U_t^{I}\theta_I)
       \left(\omega\sum_j u_j\log(1+T_j)\right).
 \tag{CI14}
\]
Here the displayed identity is understood first with an
indeterminate scalar $w$ in place of $\omega$, and then
after transporting its formal coefficients and setting
$w=\omega$.  This is why the same $\omega$ supplies the moments on
$L$: the complex scaling cancels in the function and
its directions.  No assertion that $|\iota(\lambda)|_p=1$
has been made or used.

Writing $\omega=\Omega/\iota(\delta_0)$, the number $\Omega$
is the single unit already fixed in Section~\ref{ct-section}.
Equation (CI4) implies that
$\Omega\,a\,\iota(\delta_I/\delta_0)$ is a unit as well.
This is the period adjustment that occurs when the source
changes from lattice directions to its arithmetic
trace-pairing coordinates.
The results here require $p\ge5$, $p$ split, and $p\nmid a$.
The inspected moment and local-normalization statements
explicitly impose the last condition; no claim for $p\mid a$
is needed for them.  Without it, the character in (CI9)
need not be a generator and the particular prime-to-$p$
isogeny argument above does not apply.


